\textbf{SegSieve}
\begin{lstlisting}
vector<char> segmentedSieve(long long L, long
long R) {
  // generate all primes up to sqrt(R)
  long long lim = sqrt(R);
  vector<char> mark(lim + 1, false);
  vector<long long> primes;
  for (long long i = 2; i <= lim; ++i) {
    if (!mark[i]) {
      primes.emplace_back(i);
      for (long long j = i * i; j <= lim; j += i)
        mark[j] = true;
    }
  }
  vector<char> isPrime(R - L + 1, true);
  for (long long i : primes) {
    ll sm = (L / i) * i;
    if (sm < L) sm += i;
    for (long long j = max(i * i, sm); j <= R; j
    += i)
      isPrime[j - L] = false;
  }
  if (L == 1) isPrime[0] = false;
  return isPrime;
}
\end{lstlisting}

\textbf{Phi}
\begin{lstlisting}
int phi(int n) {
  int result = n;
  for (int i = 2; i * i <= n; i++) {
    if (n % i == 0) {
      while (n % i == 0)
        n /= i;
      result -= result / i;
    }
  }
  if (n > 1)
    result -= result / n;
  return result;
}
void phi_1_to_n(){
  phi[1] = 1;
  for(int i = 2; i < N; i++) {
    if(phi[i] == 0){
      primes.push_back(i);
      phi[i] = i - 1;
    }
    for(int p : primes){
      if(i * p >= N) break;
      if(i % p == 0){
        phi[i * p] = phi[i] * p;
        break;
      }else {
        phi[i * p] = phi[i] * (p - 1);
      }
    }
  }
}
\end{lstlisting}

\textbf{Large number prime check (M.R.)}
\begin{lstlisting}
ull pow(ull a, ull t, ull mod) {
  ull r = 1;
  for (; t; t >>= 1, a = (ull)a * a % mod)
    if (t & 1) r = (ull)r * a % mod;
  return r;
}
bool isprime(ull n) {
  if (n < 2 || n % 6 % 4 != 1) return (n | 1) == 3;
  ull A[] = {2, 325, 9375, 28178, 450775,
  9780504, 1795265022},
  s = __builtin_ctzll(n - 1), d = n >> s;
  for (ull a : A) {
    ull p = pow(a % n, d, n), i = s;
    while (p != 1 && p != n - 1 && a % n && i--)
      p = (ull)p * p % n;
    if (p != n - 1 && i != s) return 0;
  }
  return 1;
}
\end{lstlisting}

\textbf{Matrix Exponentiation}
\begin{lstlisting}
vector<vector<int>>
multiply(vector<vector<int>>& a,
vector<vector<int>>& b) {
  int n = sz(a), m = sz(a[0]), p = sz(b[0]);
  vector<vector<int>> result(n, vector<int>(p,
  0));
  for (int i = 0; i < n; i++) {
    for (int k = 0; k < m; k++) {
      if (a[i][k] == 0) continue;
      for (int j = 0; j < p; j++) {
        result[i][j] = ((result[i][j] + a[i][k] *
        b[k][j]) % mod + mod) % mod;
      }
    }
  }
  return result;
}
vector<vector<int>> expo(vector<vector<int>> t,
int p) {
  int n = t.size();
  vector<vector<int>> result(n, vector<int>(n,
  0));
  for (int i = 0; i < n; i++) {
    result[i][i] = 1;
  }
  while (p > 0) {
    if (p & 1) result = multiply(result, t);
    t = multiply(t, t);
    p >>= 1;
  }
  return result;
}
\end{lstlisting}
